\section{Failure without private quadratic bounds}
\label{sec:app-rec-bounds}

We prove \cref{prop:private-bounds-needed}, including the nonmembership
and closedness facts needed for separation. This argument concerns the
scalar shared-box preordering; it does not assert closedness of the sparse
recourse certificate cone of \cref{prop:rec-dual}.

\begin{proof}[Proof of \cref{prop:private-bounds-needed}]
The arithmetic--geometric mean inequality gives
\[
  x_1^4x_2^2+x_1^2x_2^4+x_3^6
  \ge3x_1^2x_2^2x_3^2,
\]
so $q\ge0$, the objective $q(x)y^2$ is convex in $y$, and its true
minimum is zero, attained at $y=0$. Fix $r\ge3$ and put
$C=\Pre_r(\{1,2,3\})$ on the space $\Rle{x}{2r}$, which contains $q$.

First, $q\notin C$. Suppose a representation
$q=\sum_{I,l}w_Ia_{I,l}^2$ existed. Since $w_I(0)=1$ and $q(0)=0$,
every $a_{I,l}$ vanishes at the origin. Let $k$ be the smallest order of
vanishing among the nonzero $a_{I,l}$. The degree-$2k$ homogeneous part
of the representation is the sum of the squares of their degree-$k$
parts. This sum cannot vanish: a sum of real polynomial squares is zero
only if each square is zero. As $q$ is homogeneous of degree six, $k=3$.
Every $a_{I,l}$ thus has no term of degree less than three, and comparison
of degree-six parts expresses $q$ as a sum of squares of homogeneous
cubics.

Here is a direct contradiction. In any such sum, the zero coefficients
of $x_1^6$ and $x_2^6$ force the coefficients of $x_1^3$ and $x_2^3$
in every cubic to vanish. A remaining cubic has the form
\[
 a x_1^2x_2+b x_1x_2^2+c x_1^2x_3+d x_2^2x_3
 +e x_1x_3^2+f x_2x_3^2+g x_1x_2x_3+h x_3^3.
\]
The coefficient of $x_1^4x_3^2$ in the sum of squares is now $\sum c^2$,
and that of $x_2^4x_3^2$ is $\sum d^2$. Both vanish in $q$, so $c=d=0$
in every cubic. The coefficients of $x_1^2x_3^4$ and $x_2^2x_3^4$ then
give $\sum e^2=\sum f^2=0$. With these coefficients removed, the
coefficient of $x_1^2x_2^2x_3^2$ is $\sum g^2\ge0$, whereas it is $-3$
in $q$. Thus no representation exists, at any order.

Next, $C$ is closed. For each $I$ with $\abs I\le r$, let $v_I$ be the
vector of shared monomials of degree at most $r-\abs I$. A Gram
representation of an element of $C$ is
\[
 p(x)=\sum_Iw_I(x)v_I(x)^\top G_Iv_I(x),\qquad G_I\succeq0.
\]
For uniform probability measure $\omega$ on the cube, the matrix
$B_I=\int w_Iv_Iv_I^\top\,d\omega$ is positive definite: a nonzero
polynomial square has positive integral against the weight $w_I$, which
is strictly positive in the open cube. If $p_j\in C$ converges
coefficientwise, its integrals are bounded and, for any chosen Gram
representations,
\[
 \int p_j\,d\omega=\sum_I\tr(G_{I,j}B_I)
 \ge\sum_I\lambda_{\min}(B_I)\tr(G_{I,j}).
\]
All summands are nonnegative. Hence all Gram matrices are bounded, and a
subsequence of the finite collection converges to positive semidefinite
matrices giving a representation of the limit. This proves closedness.

Finite-dimensional separation now gives a linear functional $\mathcal L$
on $\Rle{x}{2r}$ such that $\mathcal L(C)\ge0$ and $\mathcal L(q)<0$. Its
mass is positive. Indeed, the unnormalized version of the proof of
\cref{lem:compact} gives $\abs{\mathcal L(x^\alpha)}\le\mathcal L(1)$ for all
$\abs\alpha\le2r$: shared generator tests bound even moments by
$\mathcal L(1)$, and Cauchy--Schwarz bounds the remaining moments. If
$\mathcal L(1)=0$, the functional vanishes on a monomial basis, contradicting
$\mathcal L(q)<0$. Scale it so that $\mathcal L(1)=1$. This normalization uses
only $r\ge3$, so it does not require $q^2$ to lie in the truncated domain.

For $R\ge0$ define on the rectangular space
\[
 L_R(a(x)+b(x)y+c(x)y^2)=a(0)+R\mathcal L(c).
\]
It has $L_R(1)=1$. For every allowed $I$ and every shared coefficient
pair $s_0,s_1$ with degree at most $r-\abs I$,
\[
 L_R\bigl(w_I(s_0+s_1y)^2\bigr)
 =w_I(0)s_0(0)^2+R\mathcal L(w_Is_1^2)\ge0.
\]
Thus (R1) holds. The two affine private localizers, with the fixed-domain
degree allowance, give
$L_R(w_Is^2(1\pm y))=w_I(0)s(0)^2\ge0$, so (R2) holds. There are no
separator equations in a single bag, and (R3) has been removed. Hence
$L_R$ is feasible for every $R$, while
$L_R(q(x)y^2)=R\mathcal L(q)\to-\infty$.
\end{proof}
